\chapter{Marfan}

\section*{Definition and Overview}
Marfan syndrome (commonly referred to simply as Marfan) is an autosomal dominant, multi-systemic connective tissue disorder characterised by variable expressivity and high penetrance. It is caused by mutations in the FBN1 gene, which encodes fibrillin-1, an essential glycoprotein that forms extracellular microfibrils and regulates the bioavailability of transforming growth factor-beta (TGF-$\beta$). Connective tissue defects in patients with Marfan syndrome manifest primarily in the cardiovascular, ocular, and musculoskeletal systems. Because of the risk of life-threatening cardiovascular complications, specifically aortic root aneurysm, dissection, and rupture, timely diagnosis and lifelong clinical surveillance are crucial.

\section*{Epidemiology}
Marfan syndrome is relatively uncommon but is one of the most frequent hereditary connective tissue disorders. It has a worldwide prevalence estimated at approximately 1 in 5,000 to 10,000 individuals, showing no predilection for geographic region, ethnicity, or sex. In Australia, thousands of individuals are affected. Around 75\% of cases are inherited from an affected parent, while the remaining 25\% arise from de novo (spontaneous) mutations. It is typically diagnosed during childhood or adolescence, although mild cases may remain undetected until adulthood.

\section*{Aetiology and Pathophysiology}
Marfan syndrome is caused by a genetic mutation affecting microfibril assembly and cytokine signalling.
\begin{itemize}
    \item \textbf{FBN1 gene mutation:} Located on chromosome 15q21.1, mutations in this gene disrupt the synthesis, secretion, or extracellular deposition of fibrillin-1.
    \item \textbf{Loss of structural integrity:} Decreased or dysfunctional fibrillin-1 leads to fragmentation of elastic fibres in the extracellular matrix of connective tissues, particularly affecting the aortic media, suspensory ligaments of the lens, and skeletal periosteum.
    \item \textbf{Dysregulated TGF-beta signalling:} Fibrillin-1 normally sequesters TGF-$\beta$ in its inactive state. Impaired fibrillin-1 binding results in excessive active TGF-$\beta$, promoting cell proliferation, inflammatory cell recruitment, and matrix metalloproteinase production, which further degrades the extracellular matrix.
    \item \textbf{Cystic medial necrosis:} In the aortic wall, loss of microfibrillar scaffolding leads to smooth muscle cell loss, elastic lamellae fragmentation, and accumulation of basophilic ground substance, making the aorta prone to progressive dilation and dissection.
\end{itemize}

\section*{Clinical Features}
The clinical manifestations of Marfan syndrome are diverse, showing significant phenotypic variability even among family members with the same mutation.
\begin{itemize}
    \item \textbf{Skeletal system:} Tall stature, disproportionately long limbs (dolichostenomelia) and fingers (arachnodactyly), joint hypermobility, pectus excavatum or carinatum, scoliosis, and a high-arched palate.
    \item \textbf{Cardiovascular system:} Aortic root dilation (sinus of Valsalva aneurysm), mitral valve prolapse (with or without regurgitation), aortic regurgitation, and a predisposition to thoracic aortic dissection.
    \item \textbf{Ocular system:} Ectopia lentis (dislocation of the lens, most commonly upward), severe myopia, early-onset cataracts, glaucoma, and an increased risk of retinal detachment.
    \item \textbf{Dermatological and pulmonary systems:} Striae atrophicae (stretch marks) in atypical locations, recurrent inguinal or incisional hernias, and an increased susceptibility to spontaneous pneumothorax.
    \item \textbf{Neurological system:} Dural ectasia (widening of the dural sac, particularly in the lumbosacral region), causing chronic lower back pain, headache, or radicular symptoms.
\end{itemize}

\section*{Differential Diagnosis}
Several other genetic conditions present with overlapping features and must be differentiated from Marfan syndrome.
\begin{itemize}
    \item \textbf{Loeys-Dietz syndrome:} Characterised by hyperelastic skin, bifid uvula, hypertelorism, and aggressive, early-onset arterial aneurysms and dissections throughout the arterial tree, caused by TGFBR1/2 or SMAD3 mutations.
    \item \textbf{Ehlers-Danlos syndrome (particularly vascular type):} Marked by joint hypermobility, tissue fragility, and arterial or bowel rupture, but lacking the classic skeletal features and ectopia lentis of Marfan syndrome.
    \item \textbf{Homocystinuria:} An autosomal recessive metabolic disorder presenting with marfanoid habitus and ectopia lentis (dislocated downward), differentiated by intellectual disability, thrombosis risk, and elevated blood methionine and homocysteine levels.
    \item \textbf{MASS phenotype:} A mild connective tissue disorder characterised by Mitral valve prolapse, Aortic diameter borderline enlargement, Skin stretch marks, and Skeletal features, without progressive aortic aneurysm or lens dislocation.
    \item \textbf{Congenital contractural arachnodactyly (Beals syndrome):} Differentiated by crumpled ears, joint contractures at birth, and lack of significant cardiovascular and ocular pathology, caused by FBN2 mutations.
\end{itemize}

\section*{When to Seek Medical Care}
Individuals with suspected features of Marfan syndrome or a family history of the condition should be referred to a geneticist or cardiologist for evaluation.

\textbf{Call 000 (or local emergency services) for:}
\begin{itemize}
    \item Sudden, tearing, severe chest, back, or abdominal pain, which is the classic presentation of acute aortic dissection.
    \item Sudden shortness of breath, chest pain worsening with respiration, or coughing up blood, which may indicate a spontaneous pneumothorax.
    \item Sudden vision loss or a sensation of a ``curtain falling'' over the eye, suggesting acute retinal detachment or lens dislocation.
\end{itemize}

\section*{Diagnosis and Investigations}
Diagnosis relies on clinical criteria combined with genetic testing, standardised by the Ghent nosology.
\begin{itemize}
    \item \textbf{Ghent criteria evaluation:} A clinical algorithm based on the presence of aortic root aneurysm, ectopia lentis, FBN1 mutation, and a systemic scoring system (which evaluates skeletal, dural, and skin features).
    \item \textbf{Echocardiography:} The primary screening and monitoring tool used to measure aortic root diameter, calculate the Z-score, and assess valvular structure and function.
    \item \textbf{Genetic testing:} DNA sequencing to identify pathogenic FBN1 mutations, helping confirm the diagnosis and facilitate screening of at-risk relatives.
    \item \textbf{Slit-lamp examination:} Performed by an ophthalmologist to evaluate for ectopia lentis (lens subluxation) and assess refractive errors.
    \item \textbf{CT or MRI angiography:} Used to obtain highly detailed, three-dimensional views of the entire aorta and other large vessels, particularly for surgical planning or monitoring areas not well visualised on echocardiography.
\end{itemize}

\section*{Management}
Management focuses on preventing aortic dissection, managing skeletal and ocular complications, and providing genetic counselling.
\begin{itemize}
    \item \textbf{Pharmacotherapy (Beta-blockers):} First-line therapy (e.g.\ metoprolol or atenolol) to reduce heart rate and myocardial contractility, lowering the haemodynamic stress on the aortic wall.
    \item \textbf{Angiotensin receptor blockers (ARBs):} Drugs such as losartan or irbesartan, which help slow aortic root dilation both by lowering blood pressure and by directly inhibiting TGF-$\beta$ signalling.
    \item \textbf{Surgical aortic replacement:} Prophylactic surgery (e.g.\ valve-sparing root replacement or composite graft replacement) is indicated when the aortic diameter reaches 50 mm (or 45 mm in those with high-risk features or family history of early dissection).
    \item \textbf{Ocular management:} Corrective lenses for myopia, surgical intervention for severe ectopia lentis or cataracts, and regular monitoring for glaucoma and retinal issues.
    \item \textbf{Lifestyle modifications:} Avoidance of strenuous contact sports, competitive athletics, and isometric exercise (e.g.\ heavy weightlifting) to prevent acute spikes in blood pressure.
\end{itemize}

\section*{Complications}
The major morbidity and mortality in Marfan syndrome are due to cardiovascular complications, though skeletal and ocular issues also cause significant disability.
\begin{itemize}
    \item \textbf{Aortic dissection and rupture:} A life-threatening tear in the inner layer of the aorta (usually Type A, involving the ascending aorta) that can lead to tamponade, ischaemia, or death.
    \item \textbf{Mitral regurgitation and heart failure:} Severe valvular incompetence leading to left ventricular enlargement and progressive congestive heart failure.
    \item \textbf{Retinal detachment and blindness:} Separation of the retina from its underlying supportive tissue, which can lead to permanent vision loss if not repaired promptly.
    \item \textbf{Severe scoliosis and chronic pain:} Spinal curvature that can compromise pulmonary function and cause debilitating muscular or skeletal pain.
    \item \textbf{Spontaneous pneumothorax:} Collapse of a lung due to rupture of subpleural blebs, causing acute respiratory distress.
\end{itemize}

\section*{Prognosis and Outcomes}
Historically, the life expectancy of patients with Marfan syndrome was significantly reduced, with death occurring in the fourth or fifth decade of life due to aortic dissection. However, with modern diagnostic methods, pharmacological therapies (beta-blockers and ARBs), and timely prophylactic aortic surgery, the life expectancy of individuals with Marfan syndrome now approaches that of the general population. Continuous, lifelong cardiac surveillance and compliance with activity restrictions are essential to maintaining these positive outcomes.

\section*{Prevention and Self-Care}
While the genetic mutation itself cannot be prevented, complications can be managed and mitigated through careful self-care and medical adherence.
\begin{itemize}
    \item \textbf{Adhere to prescribed medications:} Take daily blood pressure-lowering medications exactly as prescribed, even in the absence of symptoms.
    \item \textbf{Modify physical activities:} Opt for low-intensity, aerobic exercises (such as walking, cycling on flat terrain, or swimming) instead of high-impact or contact sports.
    \item \textbf{Attend regular follow-ups:} Keep annual appointments for echocardiograms, eye exams, and specialist consultations.
    \item \textbf{Seek genetic counselling:} Discuss inheritance risks and family planning options if you or a family member has Marfan syndrome.
    \item \textbf{Monitor for pregnancy risks:} Pregnancy significantly increases the haemodynamic strain on the aorta, requiring close maternal-foetal medicine and cardiological management.
\end{itemize}

\section*{Sources and Further Reading}
The following reputable organisations offer support, research updates, and detailed information:
\begin{itemize}
    \item The Marfan Foundation. \textit{Patient education materials, research updates, and support resources.} https://www.marfan.org/
    \item Marfan Association Queensland. \textit{Australian support network, fact sheets, and community resources.} http://www.marfanassociation.org.au/
    \item Healthdirect Australia. \textit{Overview of Marfan syndrome, diagnosis, and management guidelines.} https://www.healthdirect.gov.au/marfan-syndrome
    \item Cardiac Society of Australia and New Zealand. \textit{Clinical guidelines for the management of thoracic aortic disease.} https://www.csanz.edu.au/
    \item Mayo Clinic. \textit{Clinical overview of Marfan syndrome, symptoms, causes, and treatment options.} https://www.mayoclinic.org/diseases-conditions/marfan-syndrome/symptoms-causes/syc-20350782
\end{itemize}
